% Bookkeeping [R]. Addendum A: original AUTHOR independent second route.
% Independent stage used only the current GOAL's statements for this task.
% The named first route and sole RED were opened only after the complete
% independent save/full actual self-read and root's comparison release.

\noindent\textbf{Graft purity and the error of a measured row.}
Work on finite spaces of complete complex CAR modes, with Gaussian short-chain state,
$C_{ij}=\langle c_j^\dagger c_i\rangle$ and no anomalous covariance.
Let $V:\mathcal H_S\to\mathcal H_L$ be an isometry,
$C_S^2=C_S$, $P_S^2=P_S$, $P_L^2=P_L$ orthogonal projectors, and
$Q$ any projector on $(V\mathcal H_S)^\perp$. Set
\[
 P_0=VP_SV^\dagger+Q,\quad
 \Delta_0=V(C_S-P_S)V^\dagger,\quad D=P_0-P_L,\quad G=P_L+\Delta_0.
 \qquad K:=G^2-G=-\{D,\Delta_0\}.
\]
The mismatch $D$ includes the short--extension cross blocks; replacing
it by only its short-chain compression generally loses terms.
The graft is a state covariance only if $0\leq G\leq I$, an additional
condition not implied by a small purity defect. For Hermitian $G$ with
$\epsilon=\|K\|_{\rm op}<1/4$, define
$P_*=\mathbf1_{(1/2,\infty)}(G)$, $S=2P_*-I$, $E=G-P_*$ and
$c_\epsilon=2/(1+\sqrt{1-4\epsilon})$. Then
\[
 d_1:=\|E\|_1\leq c_\epsilon\|K\|_1.
\]
If $0\leq G\leq I$, its gauge-invariant Gaussian state obeys
\[
 d_F:=\tfrac12\|\rho_G-\rho_{P_*}\|_1
 =1-\prod_i(1-q_i)\leq d_1,\qquad q_i=\min(\lambda_i,1-\lambda_i).
\]
On this state domain $P_*$ is nearest among covariance projectors in trace norm, and its
Slater state is nearest among all pure Fock states in trace distance.
For a Hermitian one-particle row $O$, its rounding error is
\[
 \delta r=\operatorname{Tr}(OE),\quad
 |\delta r|\leq\|O\|_{\rm op}d_1,\quad
 E=S K - S E^2,\quad
 |\delta r-\operatorname{Tr}(OSK)|
 \leq\|O\|_{\rm op}\epsilon c_\epsilon^2\|K\|_1.
\]
Thus $-\operatorname{Tr}[OS\{D,\Delta_0\}]$ is a signed first-order
estimate, not by itself an error bound; the raw purity moment also
omits the spectral sign. These bounds compare to $P_*$, not an
unspecified physical continuation. For a bounded Fock observable
$\mathcal O$ the corresponding bound is $2\|\mathcal O\|_{\rm op}d_F$.
For a window $0\leq W\leq I$, put $D_W=WDW$ and
$K_W=-\{D_W,\Delta_0\}$; retain
$\tau=\|K-K_W\|_1\leq2\|(D-D_W)\Delta_0\|_1$.
Replace $\|K\|_1$ by $\|K_W\|_1+\tau$ only with a valid full-defect
operator-norm bound below $1/4$; the windowed first-order estimate
also carries $\|O\|_{\rm op}\tau$.
Common exact unitary transport preserves $d_1$, the defect norms and,
when defined, $d_F$.
A uniform observable-norm bound below a stated digit therefore persists;
one small signed row or one spatial-window reading does not.

% G1. Embedding and sign, derived directly from the supplied definitions.
% P0 and C0=P0+Delta0 are projections because their short and complement
% blocks are orthogonal. Hence Delta0^2+{P0,Delta0}-Delta0=0.
% Substitution into (PL+Delta0)^2-(PL+Delta0) gives
% {PL-P0,Delta0}=-{D,Delta0}, with D=P0-PL (the sign convention matters).
% In H_L=H_S directsum H_ext, PL=[A B;B^dag F], Delta0=[delta 0;0 0].
% Then D=[PS-A -B;-B^dag Q-F] and
% K=[{A-PS,delta} delta B;B^dag delta 0]. The external Q choice drops
% out of this anticommutator, but the cross blocks B do not. Q=0 is
% permissible as the empty-complement embedding; it does not make PL a
% disconnected vacuum. A compression of a connected vacuum is generally
% not a projector and must not be substituted for PS in this argument.
% If the exact-projector premises are relaxed, direct expansion instead
% retains V[(CS^2-CS)-(PS^2-PS)]V^dag+(PL^2-PL)-{D,Delta0}.
% The complement's defect cancels because the same Q occurs in C0,P0.
%
% Small defect is not positivity: take one short mode with PS=1, CS=0,
% and PL=|p><p| on two long modes, p=(cos theta,sin theta). The graft
% G=PL-|1><1| has eigenvalues +/-sin theta. Its defect tends to0 with
% theta, but G is not a covariance for any nonzero small theta. No Fock
% density operator or Fock trace distance is assigned to such a graft.

% G2. Distances and conventions; no covariance/Fock factor is suppressed.
% For ANY Hermitian G under the defect gap, functional calculus gives
% 2G-I=S sqrt(I+4K), hence E=2S K/(I+sqrt(I+4K)). This proves
% d1<=c_epsilon||K||_1 and ||E||_op<=epsilon c_epsilon without
% assigning a Fock state to a nonpositive or overoccupied matrix.
% For valid G diagonalize its eigenvalues lambda_i in[0,1], set
% b_i=lambda_i(1-lambda_i), q_i=min(lambda_i,1-lambda_i). Then
% q_i=2b_i/(1+sqrt(1-4b_i)), so ||K||_1<=d1 and
% b_i<=q_i<=c_epsilon b_i, with
% q_max<=epsilon c_epsilon. The gap epsilon<1/4 excludes lambda=1/2.
% To prove nearest projection, for any rank-r projection R use
% ||G-R||_1 >= Tr[(I-2R)(G-R)]=Tr G+r-2Tr(RG).
% Since 0<=R_ii<=1 and sum R_ii=r in G's eigenbasis, Tr(RG) is at
% most the sum of its r largest eigenvalues. That lower bound is
% attained by their spectral projector. Minimizing over r selects
% exactly lambda>1/2 and gives d1=sum q_i. This is not normalized
% density-matrix trace distance: if a half covariance norm is used,
% it is d1/2 and the one-particle row bound gains a factor2.
%
% The Gaussian Fock density is the product of the eigenmode Bernoulli
% states. Its largest eigenvalue is product_i(1-q_i), with eigenvector
% the Slater state of P*. For any pure |psi>, measuring its projector
% gives trace distance >=1-<psi|rho_G|psi> >=1-lambda_max(rho_G).
% The top eigenvector attains this bound (the density matrices commute).
% Thus d_F is exactly the displayed product expression and d_F<=sum q_i.
% Equivalently the FULL Fock trace norm is2d_F, not d1. For a Hermitian
% Fock observable its spectral diameter times d_F is also a valid bound.
% Nearest pure here means unrestricted particle number; if a fixed rank
% is imposed it is a separate constrained minimization. In particular,
% Tr G an integer and d1<1 imply rank P*=Tr G, so no rank change then.
% Infinite-parent claims require actual relative trace convergence and
% a common implementable Fock representation; formal sea traces do not
% establish this finite-parent statement in an infinite representation.

% G3. Exact first-order remainder and a decisive failure of a bound.
% E,S,K commute by functional calculus; P*^2=P* gives K=S E+E^2.
% Hence E=S K-S E^2 exactly, with ||S E^2||_1=Tr E^2
% <=||E||_op||E||_1<=epsilon c_epsilon^2 ||K||_1, even without
% positivity of G. Trace duality proves
% the two row inequalities. For 0<q<1/2, G=diag(q,1-q), P*=diag(0,1),
% O=diag(1,0) gives actual error q but Tr(OSK)=q(1-q)<q.
% This G is a graft of projectors: PS=|2><2|, CS=|u+><u+|,
% PL=|u-><u-|, u+/-=(sqrt(q/2),+/-sqrt(1-q/2)); V=I.
% An idle orthogonal extension makes the embedding proper if desired.
% Thus even the absolute first-order contraction need not bound error.
% Signed cancellations offer still less information. In the same graft
% family, O=diag(1,-1) gives Tr[O{D,Delta0}]=0 but Tr(OE)=2q,
% since K=-q(1-q)I. This is an algebraic graft, not a solved model row.
% Only a norm bound or an estimate WITH its controlled remainder can
% be certified from these premises. All weights/normalizations belong
% in O; for separate weighted channels sum |w_j| ||O_j|| d1_j.
% A nonlinear entropy, ratio or fitted emptying clock is not one such
% linear row. For corresponding zeros in an interval, clock error is
% <=uniform_row_error/s under a nonzero slope margin |r'|>=s there;
% existence of the compared zeros must also be secured. Integrated
% flux uses its integrated observable, whose norm may grow with time.
% For a quotient of linear rows E/L, let b_E,b_L bound the graft-to-
% rounding differences. If |L_g|>b_L, a common denominator gives
% |E_*/L_*-E_g/L_g| <=
% (b_E|L_g|+|E_g|b_L)/[|L_g|(|L_g|-b_L)]. Without that margin
% no uniform quotient bound follows from the two linear row bounds.
% Nor does nearest rounding identify an intended physical projector P_phys:
% its row error needs the additional <=||O|| ||P*-P_phys||_1 term.
% Two distinct pure projectors both have zero defect, disproving any
% identification based on purity alone; fixed Gauss/source constraints
% are not imposed by this unconstrained spectral minimization.

% G4. Window truncation is a separate error, not a purity replacement.
% T=D-WDW is Hermitian. K-K_W=-{T,Delta0}, so
% tau1<=2||T Delta0||_1 and tau_op<=2||T Delta0||_op. Also
% tau1<=2||T||_op||Delta0||_1 (or2||T||_1||Delta0||_op).
% A usable window-only estimate needs an upper bound
% epsilon_hat=||K_W||_op+tau_op<1/4 and B=||K_W||_1+tau1.
% Then d1<=c_epsilon_hat B and the row is bounded by ||O|| times this.
% For m_W=Re Tr(OS K_W), where S is the FULL graft's spectral sign,
% |delta r-m_W|<=||O||[tau1+epsilon_hat c_epsilon_hat^2 B].
% If S is approximated by S_W, additionally retain
% ||O|| ||S-S_W||_op ||K_W||_1. No commutation of K_W and S is assumed.
% If W Delta0=Delta0 W=Delta0, T Delta0=(I-W)D Delta0, so the omitted
% cross-window columns give the remainder even when the whole deviation
% lies inside the window. A small windowed defect alone is insufficient.
% If Delta0 too is replaced by Delta_W=W Delta0 W, define
% K_{W,Delta_W}=-{D_W,Delta_W}. The full-minus-truncated residual
% is K_W-K_{W,Delta_W}=-{D_W,Delta0-Delta_W}, with trace norm <=
% 2||D_W(Delta0-Delta_W)||_1. This term cannot be silently dropped.

% G5. What exact common transport proves, and what a wavefront does not.
% With G_t=U_t G U_t^dag and P*_t=U_t P* U_t^dag, spectral rounding,
% E,K,S and the covariance/Fock distances (where defined) transform
% covariantly or are invariant.
% The row is Tr[(U_t^dag O_t U_t)E]. A UNIFORM ||O_t||d1<digit proves
% a persistent absolute bound, whereas one small signed row does not.
% For a concrete valid two-mode covariance set P*=(I+sigma_x)/2,
% G=(1-q)P*+q(I-P*), E=-q sigma_x and O=(I+sigma_z)/2.
% Its initial row error is0; held h=sigma_y at t=pi/4 rotates E to
% q sigma_z and makes the row error q. All purity/distance norms stay
% fixed. This disproves persistence from a single small row, not any
% particular held-history measurement. Finite-chain returns likewise
% cannot be excluded by the purity identity.
%
% On a held h_L with [h_L,PL]=0, both Delta_t=U Delta0 U^dag AND
% D_t=U D U^dag are required in K_t=-{D_t,Delta_t}. This refers to
% the state-evolved picture. Keeping the INITIAL D and Delta0 in
% Tr[(U^dag O U){D,Delta0}] is exactly correct by trace cyclicity,
% as in section53.2's O(t) pullback; no implementation error is inferred.
% A transported window W_t=U W U^dag preserves
% the window remainder norm; a geometric moving window need not do so.
% More locally, if E=Pi_J E Pi_J+E_tail and O_t=Pi_R O_t Pi_R, then
% |delta r(t)|<=||O_t||[||Pi_R U_t Pi_J||_op^2 ||Pi_J E Pi_J||_1
%                      +||E_tail||_1]. Persistence from locality needs
% these quantities small at EVERY claimed later time; being behind a
% previously observed front alone gives no such estimate. Rounding is
% nonlocal, so localization of K_W alone does not localize E.
% Different self-consistent propagators, omitted channels or numerical
% propagation errors need their own terms; none is bounded here.

% Released comparison, within the one original author/RED process.
% SPEC section53.2's displayed purity identity agrees with G1 when its
% short state is pure quasifree, references are exact and D is the FULL
% mismatch. The phrase "whose one defect is its purity" additionally
% needs positivity; G1's independently saved example shows why.
% The parting step is "Tr(O(t)(D Delta0 + Delta0 D)) ... is the row's
% join error": it is -Tr[O(t)K], not Tr[O(t)(G-P*)]. G3's explicit
% pure-graft witness independently disproves that equality. Even the
% corrected spectral-sign first-order estimate needs its remainder.
% Section53.2(2)'s window-block wording needs G4's full products/tails.
% Section53.2(4)'s later-row condition is not a persistence theorem;
% G5 supplies the actual conditional norm statement, without ruling
% a launch, changing its operational rule or certifying an emptying fit.
%
% The released review_GRAFT_v2 agrees on exact-reference residuals,
% signed-moment limitations and the distinction from physical history.
% Its bound 2Tr(|O||K|) is independently reproduced: in each real
% eigenvalue, |lambda-round(lambda)|<=2|lambda^2-lambda|, so
% |E|<=2|K|. In O's eigenbasis |Tr OE|<=Tr(|O||E|), yielding that
% bound. Our gap gives the sharper coefficient c_epsilon<2 instead.
% Its [1/2,infinity) and our (1/2,infinity) rounding agree under the gap.
% We do not adopt its unprovided -8/81 construction or verify its code,
% fixture, filter, numerical records, gate dispositions or paper wording.
% Original RED notation: F=K, delta=epsilon, k=c_epsilon,
% RED epsilon=d1, r=epsilon*c_epsilon. Thus its r*d1 remainder is
% <= our displayed epsilon*c_epsilon^2||K||_1 by the distance bound.
% No scientific difference remains in the compared
% theorem domains. The explicit physical-target limitation and window
% Delta truncation clarification were prompted by RED/released route;
% their identities and the distinct-projector witness were rederived.
% RED's quotient bound was reproduced by the displayed common denominator.
% RED prompted the Heisenberg wording, checked here by trace cyclicity.
% No new premise or numerical acceptance has been adopted from a peer.

% SOURCES / exact access and independence.
% FULL current <internal>/GOAL_alt.md,
% including the complete owner-approved addendum A--F, actually read.
% Only its A statements supplied this task's new notation/premises.
% The current binding requirements, COMMS and rulings retain their
% full preceding-review coverage; related STATUS was not opened for A.
% Previously paired baselines, openly reused for CAR/Gaussian conventions:
% <internal>/PAGE_expectation_second_route_v1.tex
% <internal>/FIRSTLAW_rows_v1.tex
% <internal>/web40f_single_object_identity.tex
% <internal>/LIGHTCHANNEL_identities_second_route_v1.tex
% Prior paper-review intake included Page/hold numerical reports, but no
% new graft route is known to have been read. No numerical row was read
% for this task. LIGHT's first-route derivation remains unopened.
% The prior ruling exposure includes the existing graft formula and
% positivity warning in the fully read governing
% <internal>/REVIEWA_RULING_O1_D13_protocol_v1.md;
% no new blind credit is claimed for those previously available facts.
% NOT READ before the independent draft/full self-read; read only AFTER
% root's release, with the independent save safeguarded at279557ba5fd:
% <internal>/SPEC_final4b_engine.md section53.2
% <internal>/review_GRAFT_v2.tex
% <internal>/redteam_GRAFT_purity_theorem_v1.tex
% The latter two sources were read in FULL. The SPEC extraction's awk
% interval-regex stop failed, exposing section53.3--53.5 and
% section53a.0--53a.4 through EOF in addition to the authorized53.2.
% This post-draft operational overread was immediately disclosed to root;
% none supports the theorem, no independence credit is claimed for it,
% and no cited numerical-row file or implementation was opened.
% Numerical summaries in the released review/SPEC were not checked or
% adopted. LIGHT's first route and other new analytical routes stay closed.
% Only this assigned source was written. No scientific implementation,
% execution, test, gate/null readout, new numerical acceptance, foreign
% edit/copy, manuscript prose or Git mutation is part of this derivation.
